%% example-spanish-session5.tex
%% Example deck written in Spanish with the "minimalist" Beamer theme.
%% It is the original "Sesión 5" lecture ported to the theme: it shows
%% one/two/three columns, flows, tables, boxes and custom TikZ figures.
\documentclass[11pt,aspectratio=169]{beamer}

\usepackage[spanish,es-tabla]{babel}
\usetheme[footer=number]{minimalist}

\title{Sesión 5: Construcción de Software}
\subtitle{Del Diseño al Código, TDD, Calidad e Integración}
\author{Proyecto Integrador: SafeRoute}
\date{Sábado 3 de octubre de 2026}

\begin{document}

%% 1. Portada
\begin{frame}\titlepage\end{frame}

%% 2. ¿Dónde estamos?
\begin{frame}{¿Dónde estamos en el proyecto?}
Durante las sesiones anteriores construimos progresivamente una solución:

\flow[mn node/.append style={minimum width=2.4cm}]{%
  Requisitos, Modelo\\del dominio, Diseño\\OO, Diseño\\dinámico, Código}

\begin{twocol}[split=0.5]
  \begin{block}{Sesiones 2--4}
    Construimos modelos que responden:
    \begin{itemize}
      \item ¿Qué necesita el usuario?
      \item ¿Qué objetos existen?
      \item ¿Quién es responsable?
      \item ¿Cómo colaboran?
    \end{itemize}
  \end{block}
\nextcol
  \begin{block}{Sesión 5}
    Ahora debemos responder:
    \begin{center}\Large
      \textbf{¿Cómo convertimos esas decisiones en software funcional?}
    \end{center}
  \end{block}
\end{twocol}
\end{frame}

%% 3. Construcción de software
\begin{frame}{Construcción de software}
La construcción es la creación detallada de software funcional mediante
actividades como:
\begin{itemize}
  \item codificación;
  \item verificación;
  \item pruebas unitarias;
  \item integración de componentes;
  \item corrección y refactorización.
\end{itemize}

\begin{statement}
  Construir software no significa simplemente escribir código.
\end{statement}

La construcción debe respetar las decisiones tomadas durante el análisis y el diseño:

\begin{center}
  \flowline{Requisitos, Diseño, Implementación, Verificación}
\end{center}
\end{frame}

%% 4. Del diseño al código (tabla)
\begin{frame}{Del diseño orientado a objetos al código}
El diseño proporciona una guía para estructurar la implementación.

\begin{center}
  \begin{tabular}{L{4.1cm}C{1cm}L{5.3cm}}
    \toprule
    \textbf{Diseño} & & \textbf{Implementación}\\
    \midrule
    Clase         & \arr & \code{class}\\
    Atributo      & \arr & Campo o propiedad\\
    Operación     & \arr & Método\\
    Asociación    & \arr & Referencia entre objetos\\
    Multiplicidad & \arr & Colección o relación apropiada\\
    Mensaje       & \arr & Llamada a una operación\\
    \bottomrule
  \end{tabular}
\end{center}

\begin{callout}[halign=center]
  El objetivo no es traducir el UML literalmente, sino implementar las
  responsabilidades y colaboraciones que el diseño justifica.
\end{callout}
\end{frame}

%% 5. Del mensaje al método
\begin{frame}{Del mensaje del diagrama de secuencia al método}
En la sesión anterior vimos mensajes como:
\begin{center}\Large\code{envio.agregarSuministro(suministro)}\end{center}

Ese mensaje puede conducir a una operación del objeto receptor:

\flow[node distance=1.2cm,mn node/.append style={minimum width=3.3cm}]{%
  Mensaje, Operación, Implementación}

Implica que \code{Envio} debe tener una responsabilidad relacionada con la
gestión de sus suministros.

\begin{callout}[halign=center]
  La implementación debe ser consecuencia del diseño, no una actividad
  independiente de él.
\end{callout}
\end{frame}

%% 6. Construcción incremental
\begin{frame}{Construcción incremental}
En lugar de implementar todo el sistema de una sola vez:
\begin{statement}
  Construimos una pequeña parte funcional y verificable.
\end{statement}

\flow[mn node/.append style={minimum width=2.3cm,font=\footnotesize}]{%
  Historia /\\comportamiento, Prueba, Implementación, Refactorización, Integración}

Ventajas:
\begin{itemize}
  \item Reduce el tamaño de los cambios.
  \item Permite detectar errores temprano.
  \item Facilita la revisión del código.
  \item Produce incrementos funcionales.
\end{itemize}
\end{frame}

%% 7. TDD
\begin{frame}{Desarrollo guiado por pruebas: TDD}
\concept{Test-Driven Development} propone escribir primero una prueba que
exprese el comportamiento esperado.

\begin{center}
\begin{tikzpicture}[mn circle/.append style={minimum size=2.1cm}]
  \node[mn circle] (red)   {RED\\Prueba\\falla};
  \node[mn circle, right=2cm of red]   (green) {GREEN\\Prueba\\pasa};
  \node[mn circle, right=2cm of green] (ref)   {REFACTOR\\Mejorar\\diseño};
  \draw[mn arrow] (red) -- (green);
  \draw[mn arrow] (green) -- (ref);
  \draw[mn arrow, bend left=25] (ref) to (red);
\end{tikzpicture}
\end{center}

\begin{itemize}
  \item \concept{RED:} escribir una prueba que falle.
  \item \concept{GREEN:} implementar lo mínimo necesario para pasarla.
  \item \concept{REFACTOR:} mejorar el código sin cambiar el comportamiento.
\end{itemize}
\end{frame}

%% 8. Ejemplo TDD SafeRoute
\begin{frame}{TDD aplicado a SafeRoute}
\textbf{Regla de negocio:}
\begin{center}\large
  Un envío no puede superar la capacidad máxima del vehículo.
\end{center}

Primero expresamos el comportamiento mediante una prueba:

\begin{center}
\begin{tikzpicture}[mn node/.append style={minimum width=3.4cm}]
  \node[mn node] (test) {Prueba:\\envío excede capacidad};
  \node[mn node, right=1.2cm of test] (fail) {RED:\\todavía falla};
  \node[mn node, right=1.2cm of fail] (impl) {Implementación\\mínima};
  \node[mn node, below=0.8cm of impl] (pass) {GREEN:\\prueba pasa};
  \node[mn node, left=1.2cm of pass]  (ref)  {REFACTOR:\\mejorar diseño};
  \draw[mn arrow] (test) -- (fail);
  \draw[mn arrow] (fail) -- (impl);
  \draw[mn arrow] (impl) -- (pass);
  \draw[mn arrow] (pass) -- (ref);
\end{tikzpicture}
\end{center}

El ciclo permite desarrollar el comportamiento paso a paso, manteniendo una
verificación automática.
\end{frame}

%% 9. AAA (tres columnas)
\begin{frame}{Anatomía de una prueba unitaria: AAA}
Una prueba puede organizarse en tres fases:

\begin{threecol}
  \begin{block}{Arrange}
    Preparar el escenario.
    \begin{itemize}
      \item Crear objetos.
      \item Definir datos.
      \item Establecer condiciones iniciales.
    \end{itemize}
  \end{block}
\nextcol
  \begin{block}{Act}
    Ejecutar el comportamiento.
    \begin{itemize}
      \item Invocar el método.
      \item Ejecutar la operación bajo prueba.
    \end{itemize}
  \end{block}
\nextcol
  \begin{block}{Assert}
    Verificar el resultado.
    \begin{itemize}
      \item Comparar resultados.
      \item Verificar estados.
      \item Verificar errores.
    \end{itemize}
  \end{block}
\end{threecol}

\begin{callout}[halign=center]
  Una prueba debe hacer evidente qué comportamiento se está verificando y qué
  resultado se espera.
\end{callout}
\end{frame}

%% 10. Código limpio (dos columnas + tabla)
\begin{frame}{Código limpio: diseñado para ser leído}
El código se modifica y se lee constantemente. Por ello, la legibilidad es
una propiedad de calidad.

\begin{twocol}
  \begin{block}{Nombres}
    \begin{itemize}
      \item Revelar intención.
      \item Evitar abreviaturas ambiguas.
      \item Usar vocabulario del dominio.
    \end{itemize}
  \end{block}
\nextcol
  \begin{block}{Funciones}
    \begin{itemize}
      \item Responsabilidad clara.
      \item Nivel de abstracción coherente.
      \item Evitar lógica innecesariamente compleja.
    \end{itemize}
  \end{block}
\end{twocol}

\begin{center}
  \begin{tabular}{L{4.5cm}C{1cm}L{5cm}}
    \toprule
    \textbf{Menos claro} & & \textbf{Más expresivo}\\
    \midrule
    \code{calcPeso(e)} & \arr & \code{calcularPeso(envio)}\\
    \code{v}           & \arr & \code{vehiculo}\\
    \code{proc()}      & \arr & \code{validarCapacidad()}\\
    \bottomrule
  \end{tabular}
\end{center}
\end{frame}

%% 11. Code smells (tres columnas)
\begin{frame}{Code Smells: síntomas de problemas de diseño}
Un \concept{Code Smell} es un síntoma en la estructura del código que puede
indicar un problema de diseño.

\begin{threecol}
  \begin{block}{Duplicación}
    La misma lógica aparece en varios lugares.
    \par\medskip\centering\textbf{¿Dónde debería vivir esta responsabilidad?}
  \end{block}
\nextcol
  \begin{block}{Método gigante}
    Una operación concentra demasiadas tareas.
    \par\medskip\centering\textbf{¿Qué responsabilidades pueden separarse?}
  \end{block}
\nextcol
  \begin{block}{Feature Envy}
    Un método depende excesivamente de otra clase.
    \par\medskip\centering\textbf{¿Está en el objeto correcto?}
  \end{block}
\end{threecol}

\begin{callout}[halign=center]
  Los Code Smells no son necesariamente errores: son señales que justifican
  una revisión del diseño.
\end{callout}
\end{frame}

%% 12. Refactorización
\begin{frame}{Refactorización: mejorar sin cambiar el comportamiento}
La refactorización modifica la estructura interna del código sin cambiar su
comportamiento externo observable.

\flow[mn node/.append style={minimum width=3cm}]{%
  Pruebas\\verdes, Código\\funcional, Refactorización, Pruebas\\verdes}

Ejemplos:
\begin{itemize}
  \item Extraer un método.
  \item Renombrar una variable o método.
  \item Eliminar duplicación.
  \item Dividir una responsabilidad excesiva.
  \item Mejorar la estructura sin alterar el comportamiento.
\end{itemize}
\end{frame}

%% 13. Del código individual al equipo
\begin{frame}{Del código individual al trabajo en equipo}
Cuando varias personas construyen el mismo sistema necesitamos controlar los
cambios.

\begin{center}
\begin{tikzpicture}[node distance=0.5cm, mn node/.append style={minimum width=2.6cm}]
  \node[mn node] (feature) {Rama de\\trabajo};
  \node[mn node, right=of feature] (commit) {Commit};
  \node[mn node, right=of commit]  (push)   {Push};
  \node[mn node, right=of push]    (pr)     {Pull Request};
  \node[mn node, below=0.8cm of pr]  (review) {Code Review};
  \node[mn node, left=of review]     (merge)  {Integración};
  \draw[mn arrow] (feature) -- (commit);
  \draw[mn arrow] (commit) -- (push);
  \draw[mn arrow] (push) -- (pr);
  \draw[mn arrow] (pr) -- (review);
  \draw[mn arrow] (review) -- (merge);
\end{tikzpicture}
\end{center}

La integración deja de ser un evento al final del proyecto y se convierte en
una actividad controlada del proceso de desarrollo.
\end{frame}

%% 14. Branching
\begin{frame}{Estrategias de branching}
Una estrategia de branching define cómo se organiza el trabajo paralelo en un
repositorio. Para SafeRoute podemos trabajar con:

\begin{center}
\begin{tikzpicture}[mn node/.append style={minimum width=3cm,minimum height=0.65cm}]
  \node[mn node] (main) {main};
  \node[mn node, below left=0.9cm and 1.6cm of main]  (f1) {feature/calculo-peso};
  \node[mn node, below right=0.9cm and 1.6cm of main] (f2) {feature/calculo-ruta};
  \node[mn node, below=1.9cm of main] (f3) {feature/pruebas};
  \draw[mn arrow] (main) -- (f1);
  \draw[mn arrow] (main) -- (f2);
  \draw[mn arrow] (main) -- (f3);
\end{tikzpicture}
\end{center}

\flowline{Trabajo aislado, integración controlada, código compartido}

\begin{callout}[halign=center]
  Git Flow es una estrategia posible; el principio fundamental es controlar
  cómo los cambios pasan del trabajo individual a la base compartida.
\end{callout}
\end{frame}

%% 15. Commits, PR, Code Review (tres columnas)
\begin{frame}{Commits, Pull Requests y Code Review}
\begin{threecol}
  \begin{block}{Commit}
    Debe representar un cambio lógico y comprensible.
    \par\medskip Ejemplo:\par \code{feat: calcular peso total del envío}
  \end{block}
\nextcol
  \begin{block}{Pull Request}
    Propone integrar cambios en una rama compartida.
    \par\medskip Permite revisar el cambio antes de integrarlo.
  \end{block}
\nextcol
  \begin{block}{Code Review}
    Otra persona inspecciona el cambio.
    \par\medskip Se revisan comportamiento, diseño, pruebas y legibilidad.
  \end{block}
\end{threecol}

\begin{callout}[halign=center]
  El objetivo de una revisión no es encontrar culpables: es mejorar el
  producto y compartir conocimiento dentro del equipo.
\end{callout}
\end{frame}

%% 16. Construcción segura (dos columnas)
\begin{frame}{Construcción segura: validar y fallar de forma controlada}
Las entradas pueden ser incorrectas, incompletas, inesperadas o
potencialmente maliciosas.

\begin{twocol}[split=0.58]
  \concept{Principios:}
  \begin{enumerate}\small
    \item Validar las entradas en las fronteras apropiadas.
    \item Mantener invariantes del dominio.
    \item Fallar de manera explícita y controlada.
    \item No silenciar excepciones.
    \item No exponer información técnica innecesaria al usuario.
  \end{enumerate}
\nextcol
  \concept{Invariantes de SafeRoute:}
  \begin{notebox}
    \code{capacidadVehiculo > 0}\par
    \code{pesoEnvio <= capacidadVehiculo}
  \end{notebox}
\end{twocol}

\begin{callout}[halign=center]
  La seguridad también consiste en impedir que el sistema llegue a estados inválidos.
\end{callout}
\end{frame}

%% 17. SafeRoute: ejemplo completo
\begin{frame}{SafeRoute: de diseño a código y prueba}
Consideremos una responsabilidad que ya podemos derivar del diseño:
\begin{statement}Calcular el peso total de un envío\end{statement}

\flow[mn node/.append style={minimum width=2.6cm}]{DCD, Secuencia, Método, Prueba, Git}

El flujo de trabajo completo:
\begin{center}
  \flowline{Diseñar, Implementar, Probar, Refactorizar, Integrar}
\end{center}
\end{frame}

%% 18. Vertical slice
\begin{frame}{SafeRoute: una vertical slice}
En lugar de implementar todo el módulo, construiremos una pequeña
funcionalidad completa.

\flow[mn node/.append style={minimum width=2.5cm}]{%
  Historia:\\peso total, Prueba, Envio, Refactor, Pull Request}

\begin{block}{Resultado}
  Al finalizar, el equipo no solamente tendrá código: tendrá una
  funcionalidad pequeña, probada, versionada y revisada.
\end{block}
\end{frame}

%% 19. Taller
\begin{frame}{Taller: construir una funcionalidad de SafeRoute}
\textbf{Tiempo: 60 minutos}

\begin{twocol}[split=0.5]
  \begin{block}{Historia de usuario}
    Como sistema de distribución, quiero conocer el peso total de un envío
    para verificar que pueda ser transportado por un vehículo.
  \end{block}
\nextcol
  \textbf{Flujo de trabajo:}
  \begin{enumerate}\small
    \item Crear una rama \code{feature/calculo-peso-envio}.
    \item Escribir una prueba para el comportamiento esperado.
    \item Ejecutar la prueba y obtener \textbf{RED}.
    \item Implementar lo mínimo necesario.
    \item Obtener \textbf{GREEN}.
    \item Refactorizar manteniendo las pruebas verdes.
    \item Crear un commit atómico.
    \item Abrir un Pull Request.
    \item Realizar una revisión por pares.
  \end{enumerate}
\end{twocol}
\end{frame}

%% 20. Criterios de entrega
\begin{frame}{Taller: criterios de entrega}
Cada equipo deberá demostrar:
\begin{enumerate}
  \item \textbf{Repositorio:} proyecto SafeRoute disponible en Git.
  \item \textbf{Rama de trabajo:} una rama \code{feature/*} para la funcionalidad.
  \item \textbf{Código:} implementación de la responsabilidad de cálculo.
  \item \textbf{Pruebas:} al menos una prueba automatizada que demuestre el comportamiento.
  \item \textbf{TDD:} evidencia del ciclo Red \arr\ Green \arr\ Refactor.
  \item \textbf{Commit:} cambio lógico con mensaje descriptivo.
  \item \textbf{Pull Request:} solicitud de integración revisada por otro integrante.
\end{enumerate}
\end{frame}

%% 21. Preguntas de revisión (una columna angosta)
\begin{frame}{Preguntas para revisar el trabajo}
Antes de integrar el cambio:
\begin{onecol}[width=0.92]
  \begin{enumerate}\small
    \item ¿La implementación corresponde al diseño realizado anteriormente?
    \item ¿La responsabilidad está asignada al objeto adecuado?
    \item ¿La prueba describe un comportamiento real del dominio?
    \item ¿La prueba falla antes de implementar?
    \item ¿La implementación contiene solamente lo necesario?
    \item ¿El código puede leerse y entenderse fácilmente?
    \item ¿Existe duplicación o una responsabilidad excesiva?
    \item ¿La refactorización mantiene las pruebas verdes?
    \item ¿El commit representa un cambio lógico?
    \item ¿La revisión de código aporta alguna mejora concreta?
  \end{enumerate}
\end{onecol}
\end{frame}

%% 22. Cierre
\begin{frame}{Cierre: del diseño al software}
\begin{center}
  {\Large\textbf{Sesiones 3--4}}\\[0.2em]
  ¿Quién es responsable de qué?\\[0.4em]
  \large\arr\\[0.4em]
  {\Large\textbf{Sesión 5}}\\[0.2em]
  ¿Cómo convertimos esas decisiones en código funcional, probado y mantenible?\\[0.4em]
  \large\arr\\[0.4em]
  {\Large\flowline{Construir, Probar, Mejorar, Integrar}}
\end{center}

\begin{statement}
  El código es la implementación de decisiones de ingeniería, no el punto de
  partida del proyecto.
\end{statement}
\end{frame}

\end{document}
